%\documentclass[aip,reprint]{revtex4-1}
%\documentclass[prb,twocolumn,floatfix,reprint,longbibliography,superscriptaddress]{revtex4-1}
%\documentclass[prb,twocolumn,floatfix,superscriptaddress,longbibliography]{revtex4-1}
\documentclass[preprint,12pt]{elsarticle}


\usepackage{amsmath}
\usepackage{color}
\usepackage{graphicx}% Include figure files
\usepackage{dcolumn}% Align table columns on decimal point
\usepackage{bm}% bold math
%\usepackage[bookmarks=true,colorlinks,
%            citecolor=blue,linkcolor=blue, 
%           breaklinks=true]{hyperref}
%\usepackage[dvipdfm,bookmarks=true,colorlinks,citecolor=blue,linkcolor=blue, hypertex, breaklinks=true]{hyperref}
\usepackage[bookmarks=true,colorlinks,citecolor=blue,linkcolor=blue, breaklinks=true]{hyperref}
%\usepackage[bookmarks=true,colorlinks,
%            linkcolor = blue,		%%修改此处为你想要的颜色
%            anchorcolor = blue,	%%修改此处为你想要的颜色
%            urlcolor = blue,		%%修改此处为你想要的颜色
%            citecolor = blue,		%%修改此处为你想要的颜色，例如修改blue为red
%            ]{hyperref}


\usepackage{booktabs}
\usepackage{setspace}
\usepackage{subfigure}
\usepackage{enumerate}
%\makeatletter
%\newcommand{\rmnum}[1]{\romannumeral #1}
%\newcommand{\Rmnum}[1]{\expandafter\@slowromancap\romannumeral #1@}
%\makeatother


\begin{document}
	%\title{Electronic Properties of Ferromagnetic Kondo Lattice Compounds \textit{Re}Rh$_6$Ge$_4$ (\textit{Re}= Ce, Ho, Er, Tm)}
	\title{Exotic 4$f$ correlated electronic states of Ferromagnetic Kondo Lattice Compounds \textit{Re}Rh$_6$Ge$_4$ (\textit{Re}=Ce, Ho, Er, Tm)}
%\title{Complex Behavior of Lanthanides 4\textit{f} Electrons in Ferromagnetic Kondo Lattice Compounds \textit{Re}Rh$_6$Ge$_4$}

%	\author{Ping Qian}
%	\email{qianping@ustb.edu.cn}
%	\affiliation{Department of Physics, University of Science and Technology Beijing, Beijing 100083, China}
	
	
	
	
	\author{Caiqun Wang}
	\affiliation{Department of Physics, University of Science and Technology Beijing, Beijing 100083, China}
	\affiliation{Beijing Computing Center Co. Ltd, Beijing 100094, China}
	
	\author{Yu Gao}
	\affiliation{Department of Physics, Beijing Normal University, Beijing 100875, China}
	
	\author{Jun Jiang}
	\affiliation{Beijing Computing Center Co. Ltd, Beijing 100094, China}
	
	\author{Qiaoni Chen$^\dagger$}
%	\email{qiaoni@bnu.edu.cn}
	\ead{qiaoni@bnu.edu.cn}
	\affiliation{Department of Physics, Beijing Normal University, Beijing 100875, China}
	
	\author{Haiyan Lu$^\ddagger$}
%	\email{hyluphys@163.com}
	\ead{hyluphys@163.com}
	\affiliation{Institute of Materials, China Academy of Engineering Physics, Mianyang 621907, Sichuan, China}
	
	\author{Ping Qian$^\ast$}
%	\email{qianping@ustb.edu.cn}
	\ead{qianping@ustb.edu.cn}
	\affiliation{School of Mathematics and Physics, University of Science and Technology Beijing, Beijing 100083, China}
	



	
	
	\begin{abstract} 
%CeRh$_6$Ge$_4$ stands out as the first stoichiometric metallic compound with a ferromagnetic quantum critical point, thereby garnering significant attention. Ferromagnetic Kondo lattice compounds $Re$Rh$_6$Ge$_4$ ($Re$=Ce, Ho, Er, Tm) have been systematically investigated with density functional theory incorporating Coulomb interaction $U$ and spin-orbital coupling. We determined the magnetic easy axis of CeRh$_6$Ge$_4$ is within the $ab$ plane, which is in agreement with previous magnetization measurements conducted under external magnetic field and $\mu$SR experiments. We also predicted the magnetic easy axes for the other three compounds. For TmRh$_6$Ge$_4$, the magnetic easy axis aligns along the $c$ axis, thus preserving the $C_3$ rotational symmetry of the $c$ axis. Especially, there are triply degenerate nodal points along the $\Gamma-A$ direction in the band structure including spin-orbital coupling. A possible localized to itinerant crossover is revealed as 4$f$ electrons increase from CeRh$_6$Ge$_4$ to TmRh$_6$Ge$_4$. Specifically, the 4$f$ electrons of TmRh$_6$Ge$_4$ contribute to the formation of a large Fermi surface, indicating their participation in the conduction process. Conversely, the 4$f$ electrons in HoRh$_6$Ge$_4$, ErRh$_6$Ge$_4$ and CeRh$_6$Ge$_4$ remain localized, which result in smaller Fermi surfaces for these compounds. These theoretical investigations on electronic structure and magnetic properties shed deep insight into the unique nature of 4$f$ electrons, providing critical predictions for subsequent experimental studies.

		As the first metallic compound proposed to exhibit a ferromagnetic quantum critical point, CeRh$_6$Ge$_4$ has garnered significant attention. We investigate the ferromagnetic Kondo lattice compounds $Re$Rh$_6$Ge$_4$ ($Re$ = Ce, Ho, Er, Tm) systematically using density functional theory including Coulomb interaction corrections (DFT+$U$) and spin-orbit coupling. Our results indicate that the magnetic easy axis of CeRh$_6$Ge$_4$ lies within the $ab$-plane, consistent with prior magnetization measurements under external magnetic fields and $\mu\mathrm{SR}$ experiments. We also make a prediction of the magnetic easy axes for the other three compounds. For TmRh$_6$Ge$_4$, the magnetic easy axis is consistently along the $c$-axis, thereby preserving the $C_3$ rotational symmetry about the $c$-axis. Notably, in the band structure including spin-orbit coupling, triple degenerate nodes exist along the $\Gamma-A$ direction. As the number of 4\textit{f} electrons increases from CeRh$_6$Ge$_4$ to TmRh$_6$Ge$_4$, a possible transition from localized to itinerant behavior is revealed. 4\textit{f} electrons exhibit some itinerant character under conditions that are neither empty, full nor half-filled. These localized-itinerant 4\textit{f} electrons are the primary reason for the formation of the ferromagnetic ground state. Calculations of the density of states and Fermi surfaces for the compounds show that the 4\textit{f} electrons in TmRh$_6$Ge$_4$ contribute to forming a large Fermi surface, indicating their participation in the conduction process. Conversely, the 4\textit{f} electrons in HoRh$_6$Ge$_4$, ErRh$_6$Ge$_4$, and CeRh$_6$Ge$_4$ remain localized, resulting in smaller Fermi surfaces for these compounds. These theoretical investigations into the electronic structure and magnetic properties provide deep insights into the unique nature of 4\textit{f} electrons and offer key predictions for subsequent experimental studies.
	\end{abstract}
	
	\maketitle
	%\section{Introduction}
%Quantum criticality in heavy fermion materials has been a subject of ongoing experimental and theoretical research in the past twenty years. Heavy fermion materials usually contain lanthanides or actides elements which involve partially filled $f$ orbitals, so they are the typical strongly correlated electron systems. Abundant exotic phenomena have been intensively studied, such as unconventional superconductivity \cite{Ran2019_Science365-684,Jiao2020_Nature579-523--527,LiYu2021_ActaPhys.Sin.70-106-136,jiaolin:586}, non-Fermi liquid behaviors \cite{legros:cea-02086419,PhysRevLett.85.626,Daou_2008,RevModPhys.73.797,PhysRevLett.72.3262,PhysRevLett.85.626}, and exotic quantum critical phenomena \cite{Si2010_Science329-1161--1166,Liu2023_PNAS120-e2300903120,Custers2003_Nature424-524--527}. In some heavy fermion materials quantum criticality is largely driven by the competition between RKKY interaction and Kondo effect, and this mechanism falls in Landau's paradigm. However there is a large class of heavy fermion materials, on which the quantum critical point (QCP) is accompannied by the sudden breakdown of Kondo correlation. In the scenario of Kondo breakdown QCP, a jump occurs in the Fermi volume from large to small at the critical point. The Hertz-Millis-Moriya theory \cite{Hertz1976_PRB14-1165--1184,Millis1993_PRB48-7183--7196,Moriya1985} made foundation for the quantum critical phenomena of the itinerant fermions. In the ferromagnetic metallic materials, the ferromagnetic phase is either enter into other ordered phase first, or enter into the paramagentic phase through a first order phase transition \cite{Brando2016_RMP88-025006,Belitz1999_PRL82-4707--4710,PhysRevB.91.035130,PhysRevLett.105.217201,PhysRevLett.93.256404}. In the framework of Hertz-Millis-Moriya theory, two groups predicted that the ferromagnetic QCP is not stable in itinerant syste\cite{Belitz1999_PRL82-4707--4710,Chubukov2004_PRL92-147003}. However recent experiments on CeRh$_{6}$Ge$_{4}$ show clear evidence of ferromagnetic QCP under high pressure \cite{Shen_2020}.

		Over the past two decades, the origin of quantum criticality in heavy fermion materials has been a persistent focus of experimental and theoretical research. Heavy fermion materials, typically containing lanthanide or actinide elements with partially filled f-electrons, are archetypal strongly correlated electron systems. These materials often exhibit a wealth of phenomena, such as unconventional superconductivity \cite{Ran2019_Science365-684,Jiao2020_Nature579-523--527,LiYu2021_ActaPhys.Sin.70-106-136,jiaolin:586},%[1-4],
non-Fermi liquid behaviors\cite{legros:cea-02086419,PhysRevLett.85.626,Daou_2008,RevModPhys.73.797,PhysRevLett.72.3262,PhysRevLett.85.626}, %[5-9], 
and exotic quantum critical phenomena\cite{Si2010_Science329-1161--1166,Liu2023_PNAS120-e2300903120,Custers2003_Nature424-524--527}.%[10-12]. 
In some heavy fermion materials, quantum criticality is thought to be driven primarily by the competition between the RKKY interaction and the Kondo effect, a mechanism that does not exceed the Landau paradigm. However, studies have revealed a large class of heavy fermion materials where the quantum critical point (QCP) is accompanied by a sudden collapse of Kondo screening. In the Kondo breakdown scenario of the QCP, the Fermi volume undergoes a jump from a large Fermi surface to a small Fermi surface at the critical point. The Hertz-Millis-Moriya theory \cite{Hertz1976_PRB14-1165--1184,Millis1993_PRB48-7183--7196,Moriya1985}%[13-15] 
laid the groundwork for understanding quantum critical phenomena in itinerant fermion systems. In ferromagnetic metals, the ferromagnetic phase either first transitions into another ordered phase or undergoes a first-order phase transition to a paramagnetic phase \cite{Brando2016_RMP88-025006,Belitz1999_PRL82-4707--4710,PhysRevB.91.035130,PhysRevLett.105.217201,PhysRevLett.93.256404}.%[16-20]. 
Within the framework of the Hertz-Millis-Moriya theory, two research groups predicted that ferromagnetic QCPs are unstable in itinerant electron systems \cite{Belitz1999_PRL82-4707--4710,Chubukov2004_PRL92-147003}.%[17,21]. 
However, recent experiments on CeRh$_6$Ge$_4$ have provided clear evidence for the existence of a ferromagnetic QCP under high pressure \cite{Shen_2020}.%[22].

%The germanide compound CeRh$_{6}$Ge$_{4}$ is a ferromagnetic metal, and it was synthesized from the elements by Bisbuth fluxes\cite{Vosswinkel2012_ZfNatB67-1241}. Resistivity measurements under high pressure \cite{Kotegawa2019_JPSJ88-093702} first suggest there exists ferromagnetic QCP in CeRh$_{6}Ge_{4}$. Later detailed thermodynamic and transport measurements confirm the existence of the ferromagnetic QCP in CeRh$_6$Ge$_4$ under high pressure \cite{Shen_2020}. Then by silicon doping people also find ferromagentic QCP \cite{Zhang2022_PRB106-054409} under chemical pressure. Although recent theoretical work suggests that if the material is noncentrosymmetric with strong spin-orbit coupling, is quasi-one-dimensional, or is sufficiently dirty, i.e., has a short electronic mean-free path, ferromagentic QCP is possible to appear in itinerant electrons system \cite{Kirkpatrick2020_PRL124-147201}. The nature of the ferromagentic QCP on CeRh$_{6}Ge_{4}$ is still under debate. Whether it belong to the Kondo breakdown criticality \cite{Komijani2018_PRL120-157206,Shen_2020} , or there are two sequential effects near QCP as suggested by recent thermopower measurement\cite{Thomas2024_PRB109-L121105}. Recent numerical research indicate the anisotropy could lead to ferromagentci QCP \cite{Chen2022_PRB106-075114,Wang2022_SCPMA65-257211}, meanwhile both angle-resolved photoemission spectroscopy (ARPES) \cite{Wu_2021} and ultrafast \cite{Pei2021_PRB103-L180409} optical experiments suggest the anisotropic hybridization of CeRh$_{6}$Ge$_{4}$.       

The germanide CeRh$_6$Ge$_4$ is a ferromagnetic metal, synthesized from elemental precursors using a bismuth flux method \cite{Vosswinkel2012_ZfNatB67-1241}.%[23]. 
Resistivity measurements under high pressure \cite{Kotegawa2019_JPSJ88-093702} %[24] 
first hinted at the presence of a ferromagnetic QCP in CeRh$_6$Ge$_4$. Subsequent detailed thermodynamic and transport measurements confirmed the existence of a ferromagnetic QCP in CeRh$_6$Ge$_4$ under pressure \cite{Shen_2020}. %[22]. 
Later, a ferromagnetic QCP was also discovered under chemical pressure through silicon doping \cite{Zhang2022_PRB106-054409}.%[25]. 
Although recent theoretical work suggests that ferromagnetic QCPs might be possible in itinerant electron systems if the material is non-centrosymmetric with strong spin-orbit coupling, quasi-one-dimensional, or sufficiently "dirty" (i.e., having a short electron mean free path) \cite{Kirkpatrick2020_PRL124-147201}, %[26], 
the nature of the ferromagnetic QCP in CeRh$_6$Ge$_4$ remains debated. Whether it belongs to the Kondo breakdown scenario \cite{Komijani2018_PRL120-157206,Shen_2020} ,%[22,27], 
or as suggested by recent thermopower measurements hinting at two consecutive effects near the QCP \cite{Thomas2024_PRB109-L121105},%[28], 
is still under discussion. Recent numerical studies indicate that anisotropy might lead to a ferromagnetic QCP \cite{Chen2022_PRB106-075114,Wang2022_SCPMA65-257211},%[29,30], 
while angle-resolved photoemission spectroscopy (ARPES) \cite{Wu_2021} %[31] 
and ultrafast optical experiments \cite{Pei2021_PRB103-L180409} %[32] 
both demonstrate anisotropic hybridization in CeRh$_6$Ge$_4$.

%Besides CeRh$_{6}$Ge$_{4}$ other isostructure compounds $Re$Rh$_{6}$Ge$_{4}$, by substituting of $Ce$ element into other lanthanides elements have been synthesized \cite{Vosswinkel2013_Zfauac639-2623}. At the same time they all have the LiCo$_6$P$_4$ type structure \cite{buschmann1991darstellung}, and display rich exotic properties. Transport meansuments indicate that along different direction the charge carriers of LaRh$_6$Ge$_4$ are with different type \cite{Luo2023_PRB108-195146}. La-doped CeRh$_6$Ge$_4$ display the crossover from coherent Kondo lattice behaviors to Kondo impurity regime \cite{Xu2021_CPL38-087101}. DFT calculations imply there are non-trivial topological properties in the band structure of YRh$_{6}$Ge$_{4}$, LaRh$_{6}$Ge$_{4}$ and LuRh$_{6}$Ge$_{4}$ \cite{Guo2018_PRB98-045134}. Though YRh$_{6}$Ge$_{4}$, LaRh$_{6}$Ge$_{4}$ and LuRh$_{6}$Ge$_{4}$ are non-magnetic, several other $Re$Rh$_{6}$Ge$_{4}$ have local mangnetic moments. Among them HoRh$_{6}$Ge$_{4}$, ErRh$_{6}$Ge$_{4}$ and TmRh$_{6}$Ge$_{4}$ all have ferromagnetic ground state, same as CeRh$_{6}$Ge$_{4}$. In this paper we run the DFT calculations on the these four ferromagnetic germanide compounds. We first calculate the total ground state energy along different magnetic axes, and predict the easy axis of these compounds. Then we analysis the band structure, and find the nontrivial topological points in the band structure. In the end we calculate the DOS and Fermi surface, and find the evolution of the Fermi surface as the 4$f$ electrons increase. 

Besides CeRh$_6$Ge$_4$, other isostructural compounds $Re$Rh$_6$Ge$_4$ have been synthesized by replacing Ce with other lanthanide elements \cite{Vosswinkel2013_Zfauac639-2623}.%[33]. 
They all crystallize in the LiCo$_6$P$_4$-type structure  \cite{buschmann1991darstellung}, %[34] 
and exhibit a variety of exotic properties. Transport measurements indicate that carriers in LaRh$_6$Ge$_4$ have different types along different directions \cite{Luo2023_PRB108-195146}.%[35]. 
La-doped CeRh$_6$Ge$_4$ shows a crossover from coherent Kondo lattice behavior to the Kondo impurity regime \cite{Xu2021_CPL38-087101}.%[36]. 
DFT calculations suggest non-trivial topological properties in the band structures of YRh$_6$Ge$_4$, LaRh$_6$Ge$_4$, and LuRh$_6$Ge$_4$ \cite{Guo2018_PRB98-045134}.%[37].
While YRh$_6$Ge$_4$, LaRh$_6$Ge$_4$, and LuRh$_6$Ge$_4$ are non-magnetic, several other $Re$Rh$_6$Ge$_4$ compounds possess local magnetic moments. Among these, HoRh$_6$Ge$_4$, ErRh$_6$Ge$_4$, and TmRh$_6$Ge$_4$, like CeRh$_6$Ge$_4$, have a ferromagnetic ground state. In this paper, we perform DFT calculations on these four ferromagnetic germanides. We first calculate the total ground state energy along different magnetic axis directions and predict the easy axes for these compounds. We then analyze the band structures and identify non-trivial topological points within them. Subsequently, we discuss the influence of the number and nature of 4\textit{f} electrons in the lanthanide elements on the local magnetic moments of these magnetic materials. Finally, we calculate the density of states and Fermi surfaces, revealing the evolution of the Fermi surface with increasing 4\textit{f} electron count.
       	
	
	%\section{COMPUTATIONAL DETAILS}
%First-principles electronic structure calculations were performed using the fully potential (linear) augmented plane wave (LAPW)$+$local orbital (lo) method implemented in the WIEN2k package \cite{wien2k}. Exchange correlation functionals were performed using Perdew-Burke-Ernzerhof (PBE) type of generalized gradient approximation (GGA) \cite{PhysRevLett.77.3865}. A 7$\times$7$\times$11 k-points mesh is taken for the full Brillouin zone (BZ) sampling, and the maximum cut of the wave vector was adjusted to R$_{mt}$K$_{max}$ = 8.5. Murnaghan equation of state is utilized to obtain the zero temperature equilibrium structures. Once the equilibrium structures were obtained, the self-consistent calculations with an energy convergence value of $10^{-6}$ Ry and a charge convergence value of $10^{-4}$ were performed. Since the lathanides elements are heavy, the spin-orbital coupling (SOC) are considered. Due to the partially filled $4f$ orbital on the rare earth ions, the Hubbard U are also included in the calculations. So the SOC and Hubbard $U$ are both included on the rare earth ions. Ferromagnetic calculations are performed for all of the four compounds, as they all have ferromagnetic ground state. However a knotty issue in the DFT calculations of the system with $4f$ electrons is that the energy curved surface can be complex. Sometimes the calculation are difficult to converge, and sometimes they get stuck in local minima \cite{dorado2013advances}. To overcome these problems, we carefully chose the initial electronic configuration when searching for the ground state along different magnetic axes and the results are reliable.

First-principles electronic structure calculations were performed using the full-potential (linearized) augmented plane wave (FP-LAPW) plus local orbitals (lo) method as implemented in the WIEN2k software package \cite{wien2k,PRB64-195134_2001}.%[38,39].
The exchange-correlation functional was treated using the Perdew-Burke-Ernzerhof (PBE) parametrization of the generalized gradient approximation (GGA) \cite{PhysRevLett.77.3865}.%[40]. 
For full Brillouin zone (BZ) sampling, a 7×7×11 k-point mesh \cite{Brando2016_RMP88-025006}%[41] 
was employed, and the plane wave cutoff was determined by $R_{\mathrm{min}}^{\mathrm{mt}} K_{\mathrm{max}}=8.5$. The zero-temperature equilibrium structure was obtained using the Murnaghan equation of state. After obtaining the equilibrium structure, self-consistent field calculations were performed with energy convergence set to $10^{-6} \mathrm{Ry}$ and charge convergence to $10^{-4}$. Due to the presence of heavy lanthanide elements, spin-orbit coupling (SOC) was included. Given the partially filled 4\textit{f} orbitals on the rare-earth ions, a Hubbard U correction was also incorporated in the calculations. Consequently, both SOC and Hubbard U were applied on the rare-earth sites. Ferromagnetic calculations were performed for all four compounds, as they possess ferromagnetic ground states. However, a challenging aspect of DFT calculations on systems with 4\textit{f} electrons is the potentially complex energy landscape. Calculations can sometimes struggle to converge or may become trapped in local minima \cite{dorado2013advances}.%[42].
To overcome these issues, we carefully chose initial electronic configurations when searching for the ground state for different magnetic axis orientations, ensuring the reliability of the results.

	%\section{RESULTS AND ANALYSIS}
%The lattice structure of the Kondo lattice compounds $ReRh_{6}Ge_{4}$ ($Re=Ce,Ho,Er,Tm$) belong to the hexagonal crystal system. The rare-earth ions forms triangular lattice in the $ab$ plane, and the whole lattice are composed of a stack of triangular lattices along $c$ axis as displayed in Figure ~\ref{fig:struct}. The lattice constant along the $c$ axis is much shorter than it along the $a$ axis or $b$ axis as shown in Table ~\ref{tab:table1}. The $ReRh_{6}Ge_{4}$ lattice belong to the non-centrosymmetric space group $P\bar{6}m2$, and is without space-inversion symmetry. As there are two kinds of inequivalent $Rh$ irons and $Ge$ irons, marked with light (dark) blue sites and light (dark) yellows sites in Figure ~\ref{fig:struct}, the lattice is non-centrosymmetric. The first Brillioun zone is a hexagonal prism as displayed in Figure ~\ref{fig:struct} (c),  the $k_z$ direction is just along the $c$ axis, the $k_x$ direction is along the $a$ axis, and the $y$ direction is perpendicular to the $x$ direction. The red points are the high symmetry points in the Brillioun zone. The distance of $\Gamma-A$ is longer since the $c$ axis lattice constant is shorter. 

The crystal structure of the Kondo lattice compounds $Re$Rh$_6$Ge$_4$ ($Re$ = Ce, Ho, Er, Tm) is hexagonal. The rare-earth ions form a triangular lattice within the ab-plane, and the overall lattice consists of layers of this triangular lattice stacked along the c-axis, as shown in Figure 1. As shown in Table I, the lattice constant along the c-axis is significantly smaller than that along the a- or b-axes. The $Re$Rh$_6$Ge$_4$ lattice belongs to the non-centrosymmetric space group $P\bar{6}m2$, lacking spatial inversion symmetry. The presence of two inequivalent Rh ions and Ge ions, marked by light (dark) blue sites and light (dark) yellow sites in Figure 1, results in the non-centrosymmetric nature of the lattice. The first Brillouin zone is a hexagonal prism, as illustrated in ~\ref{fig:struct} (c), with the $k_z$ direction along the $c$-axis, the $k_x$ direction along the $a$-axis, and the $k_y$ direction perpendicular to $k_x$. The red points mark high-symmetry points in the Brillouin zone. The $\Gamma-A$ distance is longer due to the shorter $c$-axis lattice constant.

	\begin{figure}[htbp]
		%\centering
		%\includegraphics[width=0.48\textwidth]{figs/1.eps}
		\includegraphics[width=0.48\textwidth]{eps/jiegou.eps}
		\caption{\label{fig:struct}(a) Top view and (b) perspective view along the c-axis of the crystal structure of $Re$Rh$_6$Ge$_4$ ($Re=Ce,Ho,Er,Tm$). The red sites are rare-earth ions, and they form triangular lattice in the $ab$ plane. The blue sites are $Rh$ ions, and the yellow sites are Ge ions. (c) The first Brillioun zone of $Re$Rh$_6$Ge$_4$, and the red points are the high symmetry points.}
		\label{struct}
	\end{figure} 

%Based on the the experimental results \cite{Vosswinkel2013_Zfauac639-2623}st optimize the lattice constants of the ferromagnetic Kondo lattice compounds $ReRh$_6$Ge$_4$ ($Re$ = Ce, Ho, Er, Tm), and the results are displayed in TABLE.~\ref{tab:table1}. The lattice constants of these Kondo lattice compounds $Re$Rh$_6$Ge$_4$ are optimized by the Murnaghan equation of state \cite{murnaghan1944compressibility} in the paramagnetic phase. The optimizition start from the data of X-ray diffraction experiments on the powder samples \cite{Vosswinkel2013_Zfauac639-2623}, and our results are displayed in Table ~\ref{tab:table1}. As the atomic number of rare-earth elements increases from CeRh$_6$Ge$_4$ to TmRh$_6$Ge$_4$,  the lattice constants along $c$ axis become smaller, while the lattice constants along $a$ axis tend to become bigger. The ratio of $a/c$ becomes smaller, except from ErRh$_6$Ge$_4$ to TmRh$_6$Ge$_4$. The volume of the unit cell becomes smaller as a result of the radius of the rare-earth ions decease.

The lattice constants for these Kondo lattice compounds $Re$Rh$_6$Ge$_4$ were optimized in the paramagnetic phase using the Murnaghan equation of state \cite{murnaghan1944compressibility},%[43], 
starting from experimental X-ray diffraction data on powder samples \cite{Vosswinkel2013_Zfauac639-2623},%[33]. 
Our results are presented in Table ~\ref{tab:table1}. As the atomic number of the rare-earth element increases from CeRh$_6$Ge$_4$ to TmRh$_6$Ge$_4$, the lattice constant along the $c$-axis decreases, while the lattice constant along the $a$-axis tends to increase. The $a/c$ ratio generally decreases, except from ErRh$_6$Ge$_4$ to TmRh$_6$Ge$_4$. The unit cell volume decreases, consistent with the decreasing ionic radius of the rare-earth ions.
	\begin{table}%The best place to locate the table environment is directly after its first reference in text
		\caption{\label{tab:table1}%
			Stable lattice parameters of the rare earth germanides $ReRh_6Ge_4$ ($Re$ = Ce, Ho, Er, Tm). }
%		\begin{ruledtabular}
			\begin{tabular}{ccccc}
				\toprule
				Compound& $a$ /\r A  &$c$ /\r A &$a/c$ &Volume / \r A$^3$\\ \hline
				CeRh$_6$Ge$_4$ & 7.201 & 3.881 & 1.855 & 174.282\\
				HoRh$_6$Ge$_4$ & 7.236 & 3.840 & 1.884 & 174.124\\
				ErRh$_6$Ge$_4$ & 7.239 & 3.836 & 1.887 & 174.087\\
				TmRh$_6$Ge$_4$ & 7.238 & 3.836 & 1.886 & 174.039\\
				\bottomrule
			\end{tabular}
%		\end{ruledtabular}
	\end{table}
		
	%\subsection{Ferromagnetic Ground State and Magnetic Easy Axis}
%Magnetic measurements indicates several rare-earth germanides are with local magnetic moments \cite{Vosswinkel2013_Zfauac639-2623}, such as GdRh$_6$Ge$_4$, TbRh$_6$Ge$_4$, DyRh$_6$Ge$_4$ and YbRh$_6$Ge$_4$. Among them CeRh$_6$Ge$_4$ and the other three compounds we studied have ferromagnetic ground state. CeRh$_6$Ge$_4$ has been studied intensively, and its magnetic easy axis is within the $ab$ plane according to the magnetization measurements \cite{Shen_2020}. The magnetic Bragg peaks is not resolved in neutron diffraction of powder sample, but the coherent oscillations are observed in zero-filed $\mu SR$ measurements \cite{Shu2021_PRB104-L140411}. The $\mu SR$ measurements suggest the magnetic easy axis is along the a axis. In the following we implemented ferromagnetic DFT calculations to look for the magnetic easy axis. The spin-orbital coupling is included on $Ce$ as it's a heavy element, and Hubbard $U$ is fixed to $6$ $eV$ as previous work \cite{Wu_2021,WANG20211389}. Since lanthanides elements have unfilled $4f$ shells, the magnetism on these rare-earth germanides is mainly contributed by the lanthanides elements.     

	Magnetic measurements indicate that several rare-earth germanides possess local magnetic moments\cite{Vosswinkel2013_Zfauac639-2623}, e.g., GdRh$_6$Ge$_4$, TbRh$_6$Ge$_4$, DyRh$_6$Ge$_4$, and YbRh$_6$Ge$_4$. Among these, CeRh$_6$Ge$_4$ and the other three compounds we study exhibit a ferromagnetic ground state. CeRh$_6$Ge$_4$ has been extensively studied, and magnetization measurements show its magnetic easy axis lies within the ab-plane \cite{Shen_2020}. Neutron diffraction on powder samples could not resolve magnetic Bragg peaks, but coherent oscillations were observed in zero-field $\mu\mathrm{SR}$ measurements \cite{Shu2021_PRB104-L140411}. $\mu\mathrm{SR}$ measurements suggest the easy axis is along the $a$-axis. Next, we performed ferromagnetic DFT calculations to locate the magnetic easy axis. Given that Ce is a heavy element, spin-orbit coupling was included, and the Hubbard $U$ was fixed at 6 eV based on prior work\cite{Wu_2021,WANG20211389}. As the lanthanide elements have unfilled 4\textit{f} shells, the magnetism in these rare-earth germanides is primarily attributed to the lanthanide ions.
	\begin{figure}[htbp]
		\centering
		\includegraphics[width=0.23\textwidth]{eps/ENE-Ce.eps}
		\includegraphics[width=0.23\textwidth]{eps/ENE-Ho.eps}
		\includegraphics[width=0.23\textwidth]{eps/ENE-Er.eps}
		\includegraphics[width=0.23\textwidth]{eps/ENE-Tm.eps}
		%\label{fig:Ho-FM}
		\caption{\label{fig:ENE} The total ground state energy along different magnetic axis. (a) The ground state energy of CeRh$_6$Ge$_4$ when the magnetic axis is within the $ab$ plane. $\varphi$ is the angle with the $a$ axis, so the red line and green line represent the $a$ axis and the $b$ axis respectively. (b)-(d) The ground state energy of HoRh$_6$Ge$_4$, ErRh$_6$Ge$_4$ and TmRh$_6$Ge$_4$ when the magnetic axis is within the $ac$ plane. $\theta$ is the angle with the $c$ axis, so the blue line and the red line correspond to the $c$ axis and the $a$ axis, respectively.}
		%Total energy of (a) CeRh$_6$Ge$_4$ at different field angles $\theta$ in the $a$-$b$ plane. Total energy of (b) HoRh$_6$Ge$_4$, (c) ErRh$_6$Ge$_4$ and (d) TmRh$_6$Ge$_4$ at different field angles $\varphi$ in the $b$-$c$ plane.}		
	\end{figure} 
	
%We first calculate the ground state energy when the magnetic axis is along the lattice vector as shown in Table \ref{tab:table2}. The total ground state energy is $-1250481.63791$ $eV$ when the magnetic axis is along the $c$ axis, while the total ground state energy is $-1250481.70971$ $eV$ and $-1250481.70971$ $eV$ along the $a$ axis and $b$ axis. When the magnetic axis is along the $c$ axis the ground state energy is much higher, it indicates our calculations are in accordance with the experiments. The ground state the $ab$ plane. The $a$ axis is chosen as the $x$ axis, thus the $a$ axis correspond to $\varphi=0^{\circ}$ and the $b$ axis correspond to $\varphi=120^{\circ}$. The ground state energy of CeRh$_6$Ge$_4$ is displayed in Figure \ref{fig:ENE} (a), when the magnetic axis is within the $ab$ plane. When $\varphi=30^{\circ}$, $\varphi=90^{\circ}$ and $\varphi=150^{\circ}$, the ground state energy is the lowest. In Figure \ref{fig:ENE} (a) the ground state energy subtract the ground state energy when $\varphi=30^{\circ}$, and the value is $-1250481.72198$ $eV$. Thus our calculations imply the magnetic easy axis of CeRh$_6$Ge$_4$ is along $\varphi=30^{\circ}$, $\varphi=90^{\circ}$ and $\varphi=150^{\circ}$. 

We first calculated the ground state energy with the magnetic axis along the lattice vector directions, as shown in Table \ref{tab:table2}. For CeRh$_6$Ge$_4$, when the magnetic axis is along the $c$-axis, the total ground state energy is $-1250481.63791 \mathrm{eV}$, while along the $a$-axis and $b$-axis, the energies are $-1250481.70971 \mathrm{eV}$ and $-1250481.70971 \mathrm{eV}$, respectively. The significantly higher energy for the $c$-axis orientation confirms that our calculation aligns with experiment: the ground state is in the ab-plane. We define the $a$-axis as the $x$-axis, corresponding to $\varphi=0^{\circ}$, and the $b$-axis corresponds to $\varphi = 120^{\circ}$. The ground state energy of CeRh$_6$Ge$_4$ with the magnetic axis in the $ab$-plane is shown in Figure \ref{fig:ENE}. The lowest ground state energies occur at $\varphi=30^{\circ}$, $\varphi=90^{\circ}$, and $\varphi=150^{\circ}$. In Figure \ref{fig:ENE} (a), the ground state energy has been subtracted by the value at $\varphi = 30^{\circ}$, which is $-1250481.72198 \mathrm{eV}$. Therefore, our calculation indicates that the magnetic easy axes for CeRh$_6$Ge$_4$ are along the $\varphi=30^{\circ}$, $\varphi=90^{\circ}$, and $\varphi=150^{\circ}$ directions.
	\begin{table}%The best place to locate the table environment is directly after its first reference in text
		\caption{\label{tab:table2}%
			Ferromagnetic ground state energy of the rare earth germanides $Re$Rh$_6$Ge$_4$ ($Re$ = Ce, Ho, Er, Tm) along different axis. }
%		\begin{ruledtabular}
			\begin{tabular}{ccccc}
				\toprule
			    %Compound & $a$ axis ($eV$)   & $c$ axis ($eV$)    \\ \hline   
				%CeRh$_6$Ge$_4$ & -1250481.70971299  & -1250481.63791478 \\
				%HoRh$_6$Ge$_4$ & -1352842.01992049  & -1352843.17190210 \\
				%ErRh$_6$Ge$_4$ & -1365658.67111642  & -1365659.57456142 \\
				%TmRh$_6$Ge$_4$ & -1378779.41052308  & -1378780.40281879 \\
				%Compound & $a$ axis ($eV$)   &$c$ axis ($eV$)   &$\mu$($\mu_{B}$)  \\ \hline
                Compound & $a$ axis ($eV$)  & $b$ axis ($eV$)   & $c$ axis ($eV$)    \\ \hline   
				CeRh$_6$Ge$_4$ & -1250481.70971 & -1250481.70971 & -1250481.63791 \\
				HoRh$_6$Ge$_4$ & -1352842.01992 & -1352842.01993 & -1352843.17190 \\
				ErRh$_6$Ge$_4$ & -1365658.67112 & -1365658.67112 & -1365659.57456 \\
				TmRh$_6$Ge$_4$ & -1378779.41052 & -1378779.41052 & -1378780.40282 \\
				\bottomrule
			\end{tabular}
%		\end{ruledtabular}
	\end{table}

%The magnetization and susceptibility measurements \cite{Vosswinkel2013_Zfauac639-2623} on HoRh$_6$Ge$_4$, ErRh$_6$Ge$_4$ and TmRh$_6$Ge$_4$ indicate there are local magnetic moments on these compounds. The Curie temperature of HoRh$_6$Ge$_4$ and ErRh$_6$Ge$_4$ are all above zero, it suggests ferromagnetism on these two compounds. Both CeRh$_6$Ge$_4$ and TmRh$_6$Ge$_4$ have smaller local magnetic moment, so from the higher temperature magnetic measurements the Curie temperature is below zero. However recent experiments clarify the Curie temperature of CeRh$_6$Ge$_4$ is near $2.5K$. The local magnetic moment of TmRh$_6$Ge$_4$ is larger than CeRh$_6$Ge$_4$, thus the Curie temperature of TmRh$_6$Ge$_4$ should be higher than $2.5K$. Although the magnetic measurements are rare on HoRh$_6$Ge$_4$, ErRh$_6$Ge$_4$ and TmRh$_6$Ge$_4$, they all behave ferromagnetism. We have run the ferromagnetic DFT calculations on these three compounds. For TmRh$_6$Ge$_4$, the total ground state energy is $-1378780.40282$ $eV$ when the magnetic axis is along the $c$ axis, smaller than $-1378779.41052$ $eV$ and $-1378779.41052$ $eV$ along the $a$ axis and $b$ axis. While for HoRh$_6$Ge$_4$ and ErRh$_6$Ge$_4$, the total ground state energy along the $c$ axis is also smaller than that along the $a$ axis or $b$ axis as shown in Table \ref{tab:table2}. Thus be different from CeRh$_6$Ge$_4$, the magnetic easy axis of the other three compounds is no longer within the $ab$ plane. So we scan the ground state energy within the $ac$ plane, and the results are shown in Figure \ref{fig:ENE}. The $c$ axis is chosen as the $z$ axis, thus the $c$ axis correspond to $\theta=0^{\circ}$ and the $a$ axis correspond to $\theta=90^{\circ}$, marked with blue line and red line respectively. On HoRh$_6$Ge$_4$ and ErRh$_6$Ge$_4$ the ground state energy is lowest when $\theta=15^{\circ}$ and $\theta=165^{\circ}$ , while on TmRh$_6$Ge$_4$ the ground state energy is lowest when $\theta=0^{\circ}$. When the magnetic axis is along $\theta=15^{\circ}$ the ground state energy of HoRh$_6$Ge$_4$ and ErRh$_6$Ge$_4$ is $-1352843.19446$ $eV$ and $-1365659.64667$ $eV$ respectively, as are subtracted in Figure \ref{fig:ENE}(b) and \ref{fig:ENE}(c). The ground state energy of TmRh$_6$Ge$_4$ is $-1378780.40282$, when the magnetic axis is along the $c$ axis.         

Magnetic susceptibility and magnetization measurements  \cite{Vosswinkel2013_Zfauac639-2623} on HoRh$_6$Ge$_4$, ErRh$_6$Ge$_4$, and TmRh$_6$Ge$_4$ indicate the presence of local moments in these compounds. The Curie temperatures for HoRh$_6$Ge$_4$ and ErRh$_6$Ge$_4$ are both positive, suggesting ferromagnetism in these two compounds. CeRh$_6$Ge$_4$ and TmRh$_6$Ge$_4$ both have relatively small local moments, leading to negative Curie temperatures inferred from higher-temperature magnetic measurements. However, recent experiments clarified that the Curie temperature for CeRh$_6$Ge$_4$ is near $2.5 K$. The local moment of TmRh$_6$Ge$_4$ is larger than that of CeRh$_6$Ge$_4$, so its Curie temperature should be higher than $2.5 K$. Although magnetic measurements on HoRh$_6$Ge$_4$, ErRh$_6$Ge$_4$, and TmRh$_6$Ge$_4$ are scarce, they all exhibit ferromagnetism. We performed ferromagnetic DFT calculations for these three compounds. For TmRh$_6$Ge$_4$, the total ground state energy with the magnetic axis along the c-axis is $-1378780.40282 \mathrm{eV}$, lower than the values along the $a$-axis ($-1378779.41052 \mathrm{eV}$) and $b$-axis ($-1378779.41052 \mathrm{eV}$). Similarly, for HoRh$_6$Ge$_4$ and ErRh$_6$Ge$_4$, the total ground state energies along the $c$-axis are also lower than those along the $a$- or $b$-axes, as shown in Table \ref{tab:table2}. Thus, unlike CeRh$_6$Ge$_4$, the magnetic easy axes for the other three compounds are no longer within the $ab$-plane. Therefore, we scanned the ground state energy within the $ac$-plane; the results are shown in Figure \ref{fig:ENE}. The $c$-axis is chosen as the $z$-axis, so the $c$-axis corresponds to $\theta=0^{\circ}$, and the $a$-axis corresponds to $\theta=90^{\circ}$, marked by blue and red lines, respectively. For HoRh$_6$Ge$_4$ and ErRh$_6$Ge$_4$, the lowest ground state energies occur at $\theta=15^{\circ}$ and $\theta=165^{\circ}$, while for TmRh$_6$Ge$_4$, the minimum is at $\theta=0^{\circ}$. With the magnetic axis along $θ=15^{\circ}$, the ground state energies for HoRh$_6$Ge$_4$ and ErRh$_6$Ge$_4$ are $-1352843.19446 \mathrm{eV}$ and $-1365659.64667 \mathrm{eV}$, respectively, which are the subtracted values shown in Figures \ref{fig:ENE}(b) and \ref{fig:ENE}(c). The ground state energy for TmRh$_6$Ge$_4$ is $-1378780.40282 \mathrm{eV}$ with the magnetic axis along the $c$-axis.

	%\subsection{Band structure and potential topological properties}
%Previous DFT simulations on CeRh$_6$Ge$_4$ focus on the paramagnetic phase. We first perform the paramagnetic calculations and compare them with previous simulations, the band structure of our full potential full electron DFT calculations is consistent with these previous calculations. As the ferromagnetism play an important role in CeRh$_6$Ge$_4$, here we mainly concentrate on the ferromagnetic calculations. The ground energy along different axis indicate the easy axis is along $\varphi=30^{\circ}$, so we choose $\varphi=30^{\circ}$ as the magnetic axis. The spin-orbital coupling can not be ignored on $Ce$, and neither the Hubbard $U$. If $U$ is small the bands mainly contributed by $4f$ orbitals are quite close to the Fermi energy, and this is contraditionary with the quantum oscillation experiment \cite{WANG20211389}. In this paper we fix $U=6$ $eV$ following the previous simulations \cite{Wu_2021}. The band structure of CeRh$_6$Ge$_4$ are displayed in Figure \ref{fig:band} (a) Several bands cross the Fermi energy, so CeRh$_6$Ge$_4$ is metallic as revealed by experiments. In order to see the contribution of the $4f$ electrons, the band structure is also projected to $4f$ orbitals. The $4f$ orbitals of the $Ce^{3+}$ ions barely contribute to the bands near Fermi energy. As the electron configuration of Ce$^{3+}$ is $4f^1$, the band constitute by $4f$ orbitals would be much lower than the Fermi level.           	

Previous DFT simulations on CeRh$_6$Ge$_4$ primarily focused on the paramagnetic phase. We first performed paramagnetic calculations and compared our full-potential all-electron DFT band structures with previous simulations, finding good agreement. Since ferromagnetism plays a crucial role in CeRh$_6$Ge$_4$, we mainly focus on ferromagnetic calculations here. The ground state energy calculations for different axis directions indicate the easy axis is along $\varphi=30^{\circ}$, so we selected $\varphi=30^{\circ}$ as the magnetic axis direction for further analysis. Neither spin-orbit coupling nor the Hubbard $U$ on Ce can be neglected. A smaller U would place bands predominantly composed of 4\textit{f} orbitals too close to the Fermi level, contradicting quantum oscillation experiments \cite{WANG20211389}. %[45]. 
In this paper, following prior simulations \cite{Wu_2021}, %[31],
we fixed $U=6 \mathrm{eV}$. The band structure of CeRh$_6$Ge$_4$ is shown in Figure \ref{fig:band} (a). Several bands cross the Fermi level, confirming the metallic nature of CeRh$_6$Ge$_4$ revealed by experiments. To visualize the contribution of 4\textit{f} electrons, the band structure is projected onto the 4\textit{f} orbitals. The 4\textit{f} orbitals of $\mathrm{Ce}^{3+}$ ions contribute negligibly to the bands near the Fermi level. Given the 4$f^1$ electronic configuration of $\mathrm{Ce}^{3+}$, the bands formed by 4\textit{f} orbitals are expected to lie significantly below the Fermi level.
	\begin{figure}[htbp]
		\centering
		\includegraphics[width=0.48\textwidth]{eps/band.eps}\\
		%\includegraphics[width=0.48\textwidth]{eps/band2.eps}\\
		\caption{\label{fig:band} Band structure of (a) CeRh$_6$Ge$_4$, (b) HoRh$_6$Ge$_4$, (c) ErRh$_6$Ge$_4$ and (d) TmRh$_6$Ge$_4$ along high symmetry paths in the BZ. The calculations are carried in the ferroamgnetic phase with the SOC, and by fixed $U=6$ $eV$. These compounds are all metallic, with bands crossing the Fermi surface (red dashed lines). The blue lines are the $Re$-$4f$ orbital projected band structure. }		
	\end{figure}
	
%The band structure of HoRh$_6$Ge$_4$ and ErRh$_6$Ge$_4$ are displayed in Figure \ref{fig:band} (b) and \ref{fig:band} (c). According to the band structure, both HoRh$_6$Ge$_4$ and ErRh$_6$Ge$_4$ are metals since several bands cross the Fermi level. The electron configuration of Ho$^{3+}$ and Er$^{3+}$ are $4f^{10}$ and $4f^{11}$. The blue dots in the figures indicate the contribution from $4f$ orbitals, and there are several bands near Fermi surface are constituted by $4f$ orbitals. There are also one band cross the Fermi level with blue dots covered, and that is due to the hybridization of $4f$ orbital and itinerant orbitals. The bands mainly constituted by $4f$ orbitals in ErRh$_6$Ge$_4$ is more close to Fermi energy compared with in HoRh$_6$Ge$_4$. While in TmRh$_6$Ge$_4$ the situation is different as displayed in Figure \ref{fig:band} (d). As the number of $4f$ electrons increasing, the bands close to Fermi energy have more contribution of $4f$ orbitals, and the valley near $A$ point is lifted above Fermi level. 	

The band structures of HoRh$_6$Ge$_4$ and ErRh$_6$Ge$_4$ are displayed in Figures \ref{fig:band} (b) and \ref{fig:band} (c). According to these band structures, both HoRh$_6$Ge$_4$ and ErRh$_6$Ge$_4$ are metals, as several bands cross the Fermi level. The electronic configurations of $\mathrm{Ho}^{3+}$ and $\mathrm{Er}^{3+}$ are 4$f^{10}$ and 4$f^{11}$, respectively. The blue points indicate contributions from 4\textit{f} orbitals; several bands near the Fermi surface originate from these 4\textit{f} orbitals. Additionally, a band crossing the Fermi level is also marked with blue points, indicating hybridization between 4\textit{f} orbitals and itinerant bands. Compared to HoRh$_6$Ge$_4$, the bands primarily composed of 4\textit{f} orbitals in ErRh$_6$Ge$_4$ lie closer to the Fermi level. In TmRh$_6$Ge$_4$, however, the situation differs, as shown in Figure \ref{fig:band} (d). With the increasing number of 4\textit{f} electrons, bands near the Fermi level acquire more 4\textit{f} orbital character, and a valley near the $A$ point is pushed above the Fermi level.

%It is worth mentioning that the topological properties of the band structure. Due to the $C_{3v}$ group symmetry along the $c$ axis, and the space-inversion symmetry is not present in $ReRh_6Ge_4$, there are triply degenerate nodal points in YRh$_6$Ge$_4$, LaRh$_6$Ge$_4$ and LuRh$_6$Ge$_4$ \cite{Guo2018_PRB98-045134}. If without SOC they all have two pairs of three fold degenerate points along the $\Gamma -A$ direction. The three-fold degenerate points are generated by the crossing of a nondegenerate band and a doubly degenerate band. When SOC is included the three-fold degenerate points changes to a pair of triply degenerate nodal points. A triply degenerate nodal point is considered as intermediate state between double degenerate Weyl point and four-fold degenerate Dirac point. The $C_{3v}$ group symmetry is kept in TmRh$_6$Ge$_4$ since the magnetic axis is along the $c$ axis. There are also triply degenerate points in the band structure of TmRh$_6$Ge$_4$ along the $\Gamma -A$ direction. While in CeRh$_6$Ge$_4$, HoRh$_6$Ge$_4$ and ErRh$_6$Ge$_4$ the the $C_{3v}$ group symmetry is broken by the magnetic axis. The triply degenerate points disappear, leaving the Weyl points. 		

Worth mentioning are the topological properties of the band structures. Due to the $C_{3v}$ point group symmetry along the $c$-axis and the absence of spatial inversion symmetry in $Re$Rh$_6$Ge$_4$, triple degenerate nodes exist in YRh$_6$Ge$_4$, LaRh$_6$Ge$_4$, and LuRh$_6$Ge$_4$ \cite{Guo2018_PRB98-045134}.%[37].
Without SOC, they host two pairs of triple degenerate points along the $\Gamma-A$ direction. These triple points arise from the crossing of a non-degenerate band and a doubly degenerate band. When SOC is included, these triple points transform into a pair of triple degenerate nodes. Triple degenerate nodes are considered an intermediate state between doubly degenerate Weyl points and quadruple degenerate Dirac points. Since the magnetic easy axis of TmRh$_6$Ge$_4$ is along the $c$-axis, the $C_{3v}$ point group symmetry is preserved. In the band structure of TmRh$_6$Ge$_4$, triple degenerate points also exist along the $\Gamma-A$ direction. In contrast, for CeRh$_6$Ge$_4$, HoRh$_6$Ge$_4$, and ErRh$_6$Ge$_4$, the $C_{3v}$ symmetry is broken by the direction of the magnetic easy axis. The triple degenerate points disappear, leaving behind Weyl points.


	%\subsection{Localization and Itinerancy of 4\textit{f} Electrons}
To quantitatively characterize the gradual change in localized and itinerant behavior of the $4f$ electrons, we calculated the $4f$ charge within the Muffin-tin sphere (atomic region) and the atomic magnetic moment for the rare earth element in $Re$Rh$_6$Ge$_4$. Here $Re$ stands for all the lanthanides elements but La and Pm are excluded. The results are listed in Table \ref{tab:table3}. Along different magnetization directions, the trend in spin magnetic moment correlates with the local part of unpaired $4f$ electrons. The charge and the magnetic moment given by DFT is very little different from them in the atomic configurations. We attribute these difference to the weak overlapping of $4f$ orbitals with the orbitals from the vicinal atoms. As described by DFT means that there are very little itinerant part of the $4f$ electrons entered outside the muffin-tin sphere (the interstitial region). So we can analyzed the degree of delocalization of the $4f$ electrons in the compounds by the numerical results. Generally speak, from Ce to Eu, and then to Yb and Lu, the localized character of the 4f electrons is most pronounced when the $4f$ orbital occupation is 7 (half-filled) and 14 (fully filled).

	\begin{table}%The best place to locate the table environment is directly after its first reference in text
		\caption{\label{tab:table3}%
		The charge up/down and magnetic moments (in parentheses) from 4\textit{f} electrons in atom along different axes for the rare earth germanides $Re$Rh$_6$Ge$_4$ ($Re$ = Ce-Lu). Although SOC was considered in the calculations, band analysis shows that the contribution of $4f$ electrons near the Fermi surface is small. Therefore, the net magnetic moment from $4f$ electrons arises from the difference between spin-up and spin-down contributions, similar to the case of spin polarization.}
%		\begin{ruledtabular}
				{\fontsize{10.5pt}{5.2pt}\selectfont{
			\begin{tabular}{ccccc}
				\toprule
				Compound & & $c$-axis & $a$-axis & $ab$-plane  \\ \hline   
				CeRh$_6$Ge$_4$ & $4f^1$  & 0.38/0.15 (0.24)  & 0.58/0.14 (0.31) & 0.89/0.15 (0.78)  \\
				PrRh$_6$Ge$_4$ & $4f^3$  & 1.88/0.04 (1.89)  & 1.87/0.05 (1.87) & 1.88/0.04 (1.88) \\
				NdRh$_6$Ge$_4$ & $4f^4$  & 2.85/0.05 (2.87) & 2.85/0.05 (2.88) & 2.85/0.04 (2.88) \\
				SmRh$_6$Ge$_4$ & $4f^6$  & 4.99/0.03 (5.30) & 5.63/0.05 (4.84) & 4.85/0.04 (4.93) \\
				EuRh$_6$Ge$_4$ & $4f^7$ & 6.54/0.02 (6.63) & 6.53/0.02 (6.62) & 6.53/0.02 (6.62) \\
				GdRh$_6$Ge$_4$ & $4f^7$ & 6.87/0.04 (6.83) & 6.79/0.04 (6.87) & 6.79/0.04 (6.87) \\
				TdRh$_6$Ge$_4$ & $4f^9$ & 6.80/1.02 (5.89) & 6.80/1.02 (5.90) & 6.80/1.02 (5.90) \\
				DyRh$_6$Ge$_4$ & $4f^{10}$ & 6.77/2.38 (4.45) & 6.81/2.00 (4.89) & 6.81/2.00 (4.89) \\
				HoRh$_6$Ge$_4$ & $4f^{11}$ & 6.81/3.00 (3.86) & 6.81/3.00 (3.88) & 6.81/3.00 (3.88) \\
				ErRh$_6$Ge$_4$ & $4f^{12}$ & 6.81/4.00 (2.86) & 6.81/4.00 (2.86) & 6.81/3.91 (2.91) \\
				TmRh$_6$Ge$_4$ & $4f^{13}$ & 6.82/5.05 (1.80) & 6.80/5.01 (1.81) & 6.79/5.03 (1.78) \\
				YbRh$_6$Ge$_4$ & $4f^{14}$ & 6.74/6.74 (0.00) & 6.74/6.74 (0.00) & 6.74/6.74 (0.00) \\
		LuRh$_6$Ge$_4$ & $4f^{14}$ & 6.75/6.75 (0.00) & 6.75/6.75 (0.00) & 6.75/6.75 (0.00) \\
		\bottomrule
			\end{tabular}}}
%		\end{ruledtabular}
	\end{table}

	Before reaching half-filling of the $4f$ orbitals, as the number of unpaired $4f$ electrons increases, the electrons exhibit a trend of decreasing delocalization and increasing localization and magnetic moment (in parentheses). For instance, in Ce, the initial atomic configure is $4f^1$. Regardless of whether the magnetic axis is along the $c$ axis, $a$ axis, or the $ab$ plane, the electrons show partial occupation in both spin-up and spin-down orbitals, indicating that a small number of $f$ electrons have significant delocalized character. While as the magnetic axis evolves from the $c$ axis to the $a$ axis and the $ab$ plane, the electrons tend to become more localized, consistent with the previous calculation results. In the progression from Pr to Eu, the electrons clearly polarize. For Pr (initial atomic unpaired configure $4f^3$), the number of localized electrons is approximately 1.86; for Nd (initial $4f^4$), the number of localized electrons is about 2.86; for Sm (initial $4f^6$), the number of localized electrons is about 5.29. Clearly, the $4f$ electrons exhibit increasingly localized characteristics. Correspondingly, during the transition from half-filling to full-filling, the $4f$ orbital electron count in Gd ($4f^7$) changes from the atomic state of 7 to 6.88–6.78. However, from Tb ($4f^9\rightarrow4f^{7\uparrow}f^{2\downarrow}$), the magnetic moment (in parentheses) is 5.89, a little bigger than the net value from spin-up minus spin-down electrons in atomic configuration, due to the delocalzation mainly from spin down part. Similarly, Dy ($4f^{10}\rightarrow4f^{7\uparrow}f^{3\downarrow}$), the counterpart is 4.45-4.89 (magnetic moment); for Ho ($4f^{11}\rightarrow4f^{7\uparrow}f^{4\downarrow}$), is about 3.87 (magnetic moment); for Er ($4f^{12}\rightarrow4f^{7\uparrow}f^{5\downarrow}$), is 2.86-2.91 (magnetic moment); for Tm ($4f^{13}\rightarrow4f^{7\uparrow}f^{6\downarrow}$), is about 1.80 (magnteic moment). Finally, the $4f$ electrons of Yb and Lu have completely filled shells, resulting in zero unpaired electrons and zero magnetism, indicating complete localization.

	The localized and itinerant hybridization characteristics of the $4f$ electrons provide a good explanation for the diverse ferromagnetic ground states observed in compounds from Ce to Ho, Er, and Tm. The itinerant nature of the part-filled $4f$ orbitals in $Re$, originates from hybridization with the p-orbitals of neighboring Ge atoms. This hybridization acts as a bridge for exchange interactions between part-filled $4f$ electrons, promoting the spontaneous formation of a ferromagnetic state in these compounds. From the data, for Ce, Ho, and Er, there are differences in the distribution of localized-delocalized electrons along the $c$ axis compared to the $ab$ plane. In contrast, for Tm, $4f$ electron is localied more, due to the exist of sufficiently large number of $4f$ electrons, the few delocalized part along the $c$ axis may play a more significant role. Therefore, the complex properties exhibited by the compounds are closely related to the localized-itinerant distribution of the $4f$ electrons.


	%\subsection{Density of states and Fermi surface} 
%The density of states (DOS) of these compounds are shown in Fig.~\ref{fig:DOS}, and all the calculations are carried in the ferromagnetic case with SOC and $U=6$ $eV$. Above and below the horizontal axis are the majority and minority components of the DOS. The partial DOS of different orbitals are displayed with different colors. Purple shaded areas represent the DOS of $Rh$ $4d$ orbitals, orange lines represent the DOS of $Ge$-4$p$ orbitals, and the blue lines represent the DOS of $4f$ orbitals. The total DOS (black lines in Fig.~\ref{fig:DOS}) is non-zero at Fermi energy, since these compounds are all metallic. The DOS around Fermi energy is mainly contributed by the $4d$ orbitals of $Rh$ ions. As displayed in Figure \ref{fig:DOS}(a) there is a small blue peak near $-2$ $eV$ below the Fermi energy of CeRh$_6$Ge$_4$. It represents the occupied $4f^1$ state or the $J_{5/2}$ state considering SOC. While in HoRh$_6$Ge$_4$ and TmRh$_6$Ge$_4$ there are a few blue peaks above and close to the Fermi energy. They represent the unoccupied $4f$ orbitals. As more and more $4f$ orbitals are occupied from HoRh$_6$Ge$_4$ to TmRh$_6$Ge$_4$, one peak of the Tm$ $4f$ orbital is located at the Fermi energy. Thus the Tm$ $4f$ orbital is itinerant, and it contributes to the Fermi surface of TmRh$_6$Ge$_4$. As a result TmRh$_6$Ge$_4$ has a large Fermi surface, compared with the other three compounds on which $4f$ orbital is localized and the Fermi surface is small.     
	
The density of states (DOS) for these compounds is shown in Figure \ref{fig:DOS}. All calculations were performed for the ferromagnetic case, including SOC and $U = 6 \mathrm{eV}$. The upper and lower halves of the horizontal axis represent the majority and minority spin components of the DOS, respectively. Partial DOS from different orbitals are displayed with different colors. The purple shaded region represents the Rh-4\textit{d} orbital DOS, the orange line represents the Ge-4\textit{p} orbital DOS, and the blue line represents the 4\textit{f} orbital DOS. The total DOS (black line in Figure \ref{fig:DOS}) is non-zero at the Fermi level, confirming the metallic nature of all compounds. The DOS near the Fermi level is primarily dominated by the Rh-4\textit{d} orbitals. As shown in Figure \ref{fig:DOS}a), a small blue peak exists near $-2 \mathrm{eV}$ below the Fermi level in CeRh$_6$Ge$_4$. It represents the occupied 4f¹ state, or the $J=5/2$ state after considering SOC. In contrast, for HoRh$_6$Ge$_4$ and TmRh$_6$Ge$_4$, several blue peaks appear above but close to the Fermi level. These represent unoccupied 4\textit{f} orbitals. As more 4\textit{f} orbitals become occupied from HoRh$_6$Ge$_4$ to TmRh$_6$Ge$_4$, a peak from the Tm-4\textit{f} orbitals is situated precisely at the Fermi level. Consequently, the Tm-4\textit{f} orbitals are itinerant and contribute to the Fermi surface of TmRh$_6$Ge$_4$. As a result, TmRh$_6$Ge$_4$ possesses a large Fermi surface, unlike the other three compounds where the 4\textit{f} orbitals remain localized, leading to smaller Fermi surfaces.
	\begin{figure}[htbp]
		\centering
		%density of states\includegraphics[width=0.235\textwidth]{figs/DOS-Ce.eps}
		%\includegraphics[width=0.235\textwidth]{figs/DOS-Ho.eps}
		%\includegraphics[width=0.235\textwidth]{figs/DOS-Er.eps}
		%\includegraphics[width=0.235\textwidth]{figs/DOS-Tm.eps}
		\includegraphics[width=0.46\textwidth]{eps/DOS.eps} \\
		%\includegraphics[width=0.48\textwidth]{eps/DOS-2.eps} \\
		%\includegraphics[width=0.48\textwidth]{figs/DOS.eps} \\
		%\label{fig:Ho-FM}
		\caption{\label{fig:DOS} Density of states of (a) CeRh$_6$Ge$_4$, (b) HoRh$_6$Ge$_4$, (c) ErRh$_6$Ge$_4$ and (d) TmRh$_6$Ge$_4$. The calculations are carried in the ferromagnetic phase with the SOC, and by fixing $U=6$ $eV$. The total DOS is displayed with black line,  and the partial DOS of $Ce$-$4f$ orbitals, $Rh$-$4d$ orbitals and $Ge$-$4p$ orbitals are displayed with blue lines, shaped purple lines and orange lines. }		
	\end{figure}

%In order to shown the evolution of the Fermi surface volume, we simulate the Fermi surface of these compounds. The CeRh$_6$Ge$_4$ Fermi surface of our simulations are quite similar with previous results \cite{WANG20211389}. There are eight bands cross the Fermi energy, so the Fermi surface of CeRh$_6$Ge$_4$ is composed by eight sheets. Four of them are hole type, and four of them are electron type. The Fermi surface of HoRh$_6$Ge$_4$ and ErRh$_6$Ge$_4$ are composed by eight sheets too, and four of them are electron type. Meanwhile the Fermi surface of TmRh$_6$Ge$_4$ is composed by six sheets, and only two of them are electron type. The cross session of the Fermi surface with the $k_xk_z$ plane displayed in Figure \ref{fig:2D},  clearly shows the evolution of the Fermi surface. In TmRh$_6$Ge$_4$ the $\delta$ and $\delta'$ bands don't cross the Fermi energy, so the dark green and green lines are absent in Figure \ref{fig:2D} (d).            
  
To illustrate the evolution of the Fermi surface volume, we simulated the Fermi surfaces of these compounds. Our simulated Fermi surface for CeRh$_6$Ge$_4$ closely resembles previous results  \cite{WANG20211389}.%[45]. 
Eight bands cross the Fermi level, so the Fermi surface of CeRh$_6$Ge$_4$ consists of eight sheets. Four of these are hole-type, and four are electron-type. The Fermi surfaces of HoRh$_6$Ge$_4$ and ErRh$_6$Ge$_4$ also consist of eight sheets, with four being electron-type. Meanwhile, the Fermi surface of TmRh$_6$Ge$_4$ comprises six sheets, of which only two are electron-type. The intersections of the Fermi surfaces with the $k_xk_z$ plane are shown in Figure \ref{fig:2D}, clearly demonstrating the Fermi surface evolution. In TmRh$_6$Ge$_4$, the $\delta$ and $\delta^{\prime}$ bands do not cross the Fermi level, resulting in the absence of the dark green and green lines in Figure \ref{fig:2D} (d).
	\begin{figure}[htbp]
		\centering
		\includegraphics[width=0.48\textwidth]{eps/2D.eps} \\
		%\label{fig:Ho-FM}
		\caption{\label{fig:2D} The cross session of the $xz$ plane and the Fermi surface  of (a) CeRh$_6$Ge$_4$, (b) HoRh$_6$Ge$_4$, (c) ErRh$_6$Ge$_4$ and (d) TmRh$_6$Ge$_4$ calculated with the SOC and $U=6$ $eV$. The orange line ($\alpha$), yellow line ($\alpha'$), dark blue line ($\beta$) and blue line ($\beta'$) represent the hole type bands, while the brown line ($\gamma$), red line ($\gamma'$), dark green line ($\delta$) and green line ($\delta'$) represent the hole type bands.}		
	\end{figure}	

	%\section{SUMMARY AND OUTLOOK}
%Based on the empirical experimental evidence that the gemanides compounds $ReRh_6Ge_4$ ($Re=Ce, Ho, Er, Tm$) usually exhibit a ferromagnetic ground state, we have conducted the ferromagnetic DFT calculations on these compounds. By fixing the Coulomb interaction $U=6$ $eV$ and incorporating SOC effects, we predict the magnetic easy axis of CeRh$_6$Ge$_4$ is within $ab$ plane, with an angle of $\varphi=30^{\circ}$ with respect to the $a$ axis. For the less-explored compounds, our predictions reveal that HoRh$_6$Ge$_4$ and ErRh$_6$Ge$_4$ suggest their easy axis along the $ac$ plane, with an angle $\theta=15^{\circ}$ from the $c$ axis. In contrast, TmRh$_6$Ge$_4$ possesses a magnetic easy axis aligned along the $c$ axis, which preserves the $C_{3v}$ point group symmetry. Intriguingly, within the $\Gamma-A$ direction, we observe the presence of triply degenerate nodal points in their band structures. The magnetic axes of the other three compounds, as a result of their unique magnetic properties, break the $C_{3v}$ symmetry, even though the nodal points are absent. Instead, Weyl points emerge along the $\Gamma-A$ direction. The density of states (DOS) analysis further elucidates the almost localized 4$f$ states, and enhancing itinerancy of 4$f$ states from HoRh$_6$Ge$_4$ to TmRh$_6$Ge$_4$. TmRh$_6$Ge$_4$ boasts a substantial Fermi surface due to the active participation of the 4$f$ electrons in its Fermi surface formation.

	Based on experimental evidence that the germanides $Re$Rh$_6$Ge$_4$ (Re = Ce, Ho, Er, Tm) typically exhibit a ferromagnetic ground state, we performed ferromagnetic DFT calculations on these compounds. By fixing the Coulomb interaction $U=6 \mathrm{eV}$ and including SOC effects, we predict that the magnetic easy axis of CeRh$_6$Ge$_4$ lies within the ab-plane, at an angle $\varphi=30^{\circ}$ relative to the $a$-axis. For the less-studied compounds, our predictions indicate that the easy axes of HoRh$_6$Ge$_4$ and ErRh$_6$Ge$_4$ lie within the $ac$-plane, at an angle $\theta=15^{\circ}$ relative to the $c$-axis. In contrast, the magnetic easy axis of TmRh$_6$Ge$_4$ is along the $c$-axis, preserving the $C_{3v}$ point group symmetry. Interestingly, within the $\Gamma-A$ direction, we observe triple degenerate nodes in their band structures. The other three compounds, due to their unique magnetic orientations, have magnetic axes that break the $C_{3v}$ symmetry; thus, while nodes are absent, Weyl points appear along the $\Gamma-A$ direction. Analysis of the density of states further elucidates that the 4\textit{f} states are largely localized, with itinerancy increasing from HoRh$_6$Ge$_4$ to TmRh$_6$Ge$_4$. Because the 4\textit{f} electrons actively participate in forming the Fermi surface, TmRh$_6$Ge$_4$ hosts a large Fermi surface.

%Recent thermoelectric measurements \cite{Thomas2024_PRB109-L121105} reported an intriguing phenomenon in CeRh$_6$Ge$_4$, where the onset of orbital-selective hybridization within the ferromagnetic phase at around $0.7$ $GPa$ leads to a discernible alteration in the Fermi surface geometry. Despite this, quantum oscillation data indicate that the extremal orbits on the Fermi surface remain unaltered across the critical temperature $T_c$. This discrepancy underscores the need for further numerical investigations on CeRh$_6$Ge$_4$ under pressure to delve deeper into the physics near the quantum critical point (QCP). Meanwhile, experimental studies on the isostructural counterparts, HoRh$_6$Ge$_4$, ErRh$_6$Ge$_4$ and TmRh$_6$Ge$_4$, remain relatively sparse. Our results shall establish a connection between magnetism and electronic structure, facilitating subsequent theoretical and experimental research. 

	Recent thermoelectric measurements reported an intriguing phenomenon in CeRh$_6$Ge$_4$: the onset of orbital-selective hybridization within the ferromagnetic phase at about $0.7 \mathrm{GPa}$ leads to a perceptible change in the Fermi surface geometry. Nonetheless, quantum oscillation data indicate that the extremal orbits on the Fermi surface remain unchanged across the critical temperature $T_c$. This discrepancy underscores the need for further numerical studies on CeRh$_6$Ge$_4$ under pressure to delve deeper into the physics near the quantum critical point. Simultaneously, experimental investigations on the isostructural counterparts HoRh$_6$Ge$_4$, ErRh$_6$Ge$_4$, and TmRh$_6$Ge$_4$ remain relatively sparse. Our results will establish a connection between magnetic properties and electronic structure, facilitating subsequent theoretical and experimental efforts.

%\begin{acknowledgments}

This work is supported by the  National Key Research and Development Program of China (under Grant No. 2024YFA1408600), the National Natural Science Foundation of China (under Grant No. 12474241), Presidential Foundation of CAEP (under Grant No. YZJJZQ2024014), National Key Research and Development Program of China (under Grant No. 2022YFA1402201).  We are grateful for the fruitful conversations with Pengjie Guo, and we thank Kai Liu for helpful discussions. We acknowledge National Super Computer Center in Tianjin for computing time.
%This work was supported by the National Natural Science Foundation of China (Grant Nos. 11974048, 11974049, 2022YFA1402201) and a project from the China Academy of Engineering Physics (No. TCGH0710). We thank Kai Liu and Pengjie Guo for helpful discussions. We acknowledge the National Supercomputing Center in Tianjin for providing computational time.
%\end{acknowledgments} 

%\bibliographystyle{apsrev4-1}
\bibliographystyle{elsarticle-num} 
\bibliography{RE}

\end{document}
